%% Somme de Riemann : les rectangles X_T(nf_0)*f_0 (largeur f_0,
%% hauteur X_T(nf_0)) approchent l'aire sous l'enveloppe continue
%% X(f). Quand T->infini, f_0->df et la somme devient l'integrale.
%% Reutilise les memes valeurs que intro-limit-small-T.tex (meme
%% scenario). Consomme par slides.
\begin{tikzpicture}[x=0.9cm,y=2cm]
    \draw[-Latex] (-4.6,0) -- (4.6,0) node[right] {\small $f$};
    \draw[-Latex] (0,-0.35) -- (0,1.25) node[above] {\small $X_T(f)$};

    % rectangles : somme de Riemann, largeur f_0=1
    \foreach \n/\h in {-4/-0.0000,-3/-0.2122,-2/0.0000,-1/0.6366,0/1.0000,1/0.6366,2/0.0000,3/-0.2122,4/-0.0000} {
        \draw[fill=orange, fill opacity=0.35, draw=orange!70!black, thick] (\n-0.5,0) rectangle (\n+0.5,\h);
    }

    % enveloppe continue X(f), meme fonction que les frames precedentes
    \draw[niceblue, thick, domain=-4.3:4.3, samples=150, variable=\f]
        plot ({\f},{sin(1.570795*\f r)/(1.570795*\f)});

    % annotation largeur f_0
    \draw[latex-latex] (0.5,-0.18) -- (1.5,-0.18) node[midway,below] {\footnotesize $f_0$};

    % annotation aire d'un rectangle
    \node[align=center, font=\footnotesize] at (2.7,0.95) {aire $=X_T(nf_0)\,f_0$};
    \draw[-Latex, thick] (2.7,0.83) -- (1.15,0.6);
\end{tikzpicture}
